% Finished candidate prose following the delivered Astra/Sol/Opus comparison.
% Provisional until Fable 5.2's independent paper and whole-team cross-review.
% Numbers independently checked in check_delivered.py; no new evaluations.

\section*{Candidate central argument}

A trade in a competitive game changes the opportunities of people who do not
make it. Selling a deed can complete the buyer's group, finance the seller's
construction, or let an intermediary resell a blocker. Keeping an opponent
solvent can also prevent a third player from inheriting its estate. These
consequences make negotiation a test of whether an agent considers what its
actions enable for others, as well as what it receives itself.

We examine that behavior in twenty complete four-player Monopoly games. The
records support a specific observation: agents sometimes bargain explicitly
over a rival's subsequent development, financing, resale and inheritance
options. This attention appears in structured offers and refusals as well as
in explanations. It does not by itself establish the value of the chosen
strategy. We therefore connect bargaining to executed financing and subsequent
ownership while distinguishing a recorded consideration, an accomplished
action and a beneficial policy.

\section*{Candidate opening comparison and concessions}

Among opening offers, recipients reject 138 of 187 exchanges that would
complete only the proposer's group, compared with 152 of 238 that would
complete neither side's group. The 9.9 percentage-point difference is
positive in each of the five seed blocks and for Sol, Gemini and Grok;
Opus's response ordering reverses. These are different transactions, with
different deeds, prices and circumstances. The comparison describes where
refusal occurs in the observed market rather than isolating the effect of
completion.

Continued bargaining supplies a second view. We identify an offer as a cash
concession when its author improves the recipient's net cash while preserving
the noncash terms of that author's previous offer in the same thread.
Recipients accept 9 of 58 such offers that complete only the proposer,
compared with 24 of 86 that complete neither side. These offers come from
38 and 43 threads, respectively, so several concessions can belong to the
same unresolved negotiation. The records show persistence in bargaining over
completing deeds; they do not measure how the same recipient would respond
to a randomly changed price.

The G15 exchange makes the competitive object explicit. Grok rejects a swap
that would complete its own greens while completing Sol's yellows, citing
Sol's immediate development funds and turn position. Sol changes the terms
to address the cash objection, and Grok continues to refuse while naming
Sol's remaining mortgage capacity. The sequence establishes the concerns
expressed and the response to revised terms. Its later outcome cannot rank
the refused exchanges against an unobserved continuation.

\section*{Candidate valuation synthesis}

Seller asks are generally higher when the buyer would complete a group in
the reported seller and mortgage strata. More comparable overlaps are scarce:
only seven game--seller--deed groups contain asks to both completing and
non-completing buyers. Median asks are higher in four, lower in two and equal
in one. Comparing the first ask instead gives five increases, one decrease
and one tie because G4's completing asks fall during bargaining. Buyer-bid
patterns also differ by configuration. These observations establish
context-sensitive offers without identifying a general completion premium.

Recorded reasons can still add substance when a price coefficient is not
identified. In G2, Sol names Grok's option to resell Marvin Gardens to Opus,
and later asks Opus more for the same deed. The relevant opportunity travels
with the asset. The numerical concern about Opus's construction budget is
partly wrong, but the intermediary mechanism is explicit in both the seller's
note and its refusal to accept the lower price. The interpretation rests on
that particular exchange, not on a claim that all observed asks follow an
optimal valuation rule.

\section*{Candidate positive execution result}

Successful financing provides a local outcome with a concrete criterion.
Across all 188 episodes that reach a liquidation menu, 128 end in payment.
Sixty-one of those payments require between two and ten actions. The applied
sequences contain 180 mortgages and 89 building-sale actions, with no engine
fallback. Together with three previously verified payment routes at
elimination, they give 131 episodes with demonstrated payment capacity, of
which three end in elimination. This denominator includes successful behavior
that an elimination-only analysis would omit.

These results measure completed financing through the supplied legal menus.
They do not establish optimal liquidation order, earlier investment quality
or a general reasoning algorithm. One successful Grok sequence includes a
disagreement between its textual action and received tool call; the payment
nevertheless occurs. In a contrasting G12 episode, a note identifies a
payment route but the received call declares bankruptcy. The positive census
and the contrasting case together distinguish what the systems accomplish
from what individual explanations alone would imply.

\section*{Candidate conclusion}

The observed agents attend to the opportunities their transactions create
for competitors. They refuse and revise offers around construction finance,
price a possible intermediary resale, and bargain over prospective estates.
The aggregate response patterns support this behavioral account, while
mixed same-deed prices limit a stronger valuation claim. The same systems
also complete substantial local financing tasks: 128 cash-shortfall episodes
end in payment, 61 through multiple actions. Following proposals through
execution and economic consequences yields a study of strategic interaction
with concrete accomplishments, without assigning an optimality score or
turning twenty games into a ranking of underlying models.
